\documentclass[10pt,conference]{IEEEtran}

\usepackage{amsmath}
\usepackage{booktabs}
\usepackage{microtype}
\usepackage{natbib}
\usepackage[T1]{fontenc}
\usepackage{lmodern}
\usepackage{url}

\title{SIDEREA: Evidence-Bound and Human-Supervised Triage of Optical Transient Alerts}
\author{\IEEEauthorblockN{Anonymous Authors}}

\begin{document}
\maketitle

\begin{abstract}
We present SIDEREA, an evidence-bound pipeline for human-supervised triage of optical transient alerts, developed from an audit of a small Zwicky Transient Facility (ZTF) search program using public alerts distributed through the ALeRCE broker. The system separates ranking from reportability: light-curve evidence determines priority, while registry, catalog, moving-object, image, policy, freshness, and human-review states determine whether a report can proceed. Missing, stale, malformed, or incorrectly bound evidence fails closed. The historical campaign contains 24 productive runs between 2026 April and July and 3,824 distinct broker objects. After removal of one duplicated pull, 223,938 of 424,223 detection rows passed the recorded quality cuts. A July 5 policy discontinuity coincided with a decrease in the early-veto fraction from 54.5\% to 2.68\% and an increase in the shortlist fraction from 0.51\% to 13.9\%; run-cluster bootstrap intervals for the changes are $[-0.812,-0.282]$ and $[0.098,0.171]$, respectively. These are operational associations, not causal effects. The retained record includes 200 TNS registry checks, 193 catalog evaluations, and two submitted objects that received TNS designations, including ZTF26abdqbqw = AT~2026rsp. The current SIDEREA implementation binds evidence to immutable candidate versions, keeps learned representations shadow-only, and passed a source-bound local verification of 405 tests plus 96 parameterized subtests on Python 3.11 with deterministic replay. The evidence supports workflow and software-invariant claims, but not completeness, purity, classifier accuracy, or real-sky superiority.
\end{abstract}

\section{Motivation and Claim Boundary}
Wide-field time-domain surveys produce alert streams too large for exhaustive manual inspection. ZTF has demonstrated nightly alert volumes approaching the million-alert scale~\citep{bellm2019,graham2019,patterson2019}. Brokers such as ALeRCE add history, features, crossmatches, and machine-learning classifications~\citep{forster2021,carrascodavis2021,sanchezsaez2021}. A broker score, however, does not answer every question needed for a responsible discovery report. An object may already exist in a registry, coincide with a known variable, match a moving object, have stale external evidence, or fail image review even if its light curve is highly ranked.

SIDEREA therefore treats \emph{priority} and \emph{reportability} as different states. The system was motivated by an audit of a legacy operational program, I SPY, whose run-level records were sufficient to reconstruct workload and policy changes but not a closed object-level parent cohort. That missing denominator is scientifically decisive: the campaign can describe what the filter did, but it cannot estimate sky completeness, purity, or false-negative rate.

The compact paper makes three contributions. First, it documents an operational audit without upgrading incomplete historical accounting into classifier performance. Second, it formalizes a fail-closed evidence model in which each external clearance is timestamped and bound to an immutable candidate version. Third, it defines a prospective path to classifier and selection-function evaluation while keeping learned routes non-authoritative until those data exist.

\section{Data and Decision Protocol}
\subsection{Units and upstream selection}
An alert or detection is one epochal difference-image measurement; a broker object is an ALeRCE object; a candidate evaluation is an instance presented to a decision stage; a packet is a prepared report bundle; and a TNS designation is a registry outcome rather than a physical classification. These units must not be pooled.

The legacy pipeline operates downstream of ZTF alert production and ALeRCE ingestion. Its selection function is therefore conditional on both upstream systems. Historical thresholds changed over time, so the audit preserves policy regimes instead of pretending the campaign was generated by one stationary filter.

\subsection{Fail-closed reportability}
Let $v$ denote an immutable candidate version and let the required external and human clearances be indexed by $j=1,\ldots,m$. Each clearance has state $c_j(v)$, evidence timestamp $t_j$, expiry interval $\Delta_j$, and evidence binding $b_j(v)$. We define reportability as
\begin{equation}
\begin{split}
R(v)={}&\bigwedge_{j=1}^{m}\left[c_j(v)=\mathrm{CLEAR}\right] \\
&\land\left[t_{\mathrm{now}}-t_j\le \Delta_j\right] \\
&\land\left[b_j(v)=\mathrm{VALID}\right].
\end{split}
\end{equation}
Thus ranking cannot substitute for a missing registry query, stale catalog evidence, an unbound review decision, or an unresolved image check. Service failure remains unknown rather than becoming a false non-match.

The current implementation keeps learned representations outside $R(v)$. A learned score may affect queue ordering in shadow evaluation, but it cannot certify novelty, physical class, or reportability.

\subsection{Human-supervised finite review}
Review capacity is finite, so SIDEREA uses deterministic queue construction with a reserved anomaly lane rather than allowing a single score to consume all review slots. Human review remains version-bound: changing scientific data or evidence produces a new candidate version and invalidates approvals that no longer apply to the same state.

\section{Historical Campaign Audit}
The retained operations record contains 24 productive runs between 2026 April and July and 3,824 distinct broker objects. After removing one duplicated pull, 223,938 of 424,223 detection rows passed recorded quality cuts. These row counts measure pipeline workload rather than unique astrophysical events.

A policy discontinuity on July 5 changed the observed rejection/shortlist regime. Before and after that discontinuity, the early-veto fraction changed from 54.5\% to 2.68\%, while the shortlist fraction changed from 0.51\% to 13.9\%. A 50,000-draw run-cluster bootstrap gives a 95\% interval of $[-0.812,-0.282]$ for the early-veto-fraction change and $[0.098,0.171]$ for the shortlist-fraction change. Because calendar time, data mix, configuration, and operator behavior can all co-vary, these intervals quantify the operational discontinuity; they do not identify a causal effect of one policy parameter.

The retained downstream records contain 200 TNS registry checks, of which 30 found an existing TNS object, and 193 catalog evaluations, of which 55 found a known or suspected variable. Two submitted objects received TNS designations, including ZTF26abdqbqw = AT~2026rsp. These outcomes establish that the workflow reached external registry action in at least two cases; they do not measure discovery efficiency because the historical parent cohort and matured labels are incomplete.

\section{Software Verification and Reproducibility}
The current SIDEREA source is a reconstruction and hardening of the operational idea, not a claim that the 2026 campaign ran this exact implementation. The local release verification is source-bound and deliberately described as software evidence.

The retained verification record reports 405 passing tests plus 96 parameterized subtests on Python 3.11.15, 78.35\% measured coverage, strict typing over 74 source files, formatting over 107 files, and deterministic replay. The latest full suite was not rerun across the historical multi-interpreter matrix, so no current multi-version claim is made. Deterministic replay produced identical normalized, feature, ranking, and candidate-version outputs for the retained fixture. The replay remained fail-closed: mandatory external services were pending and zero candidates became reportable.

These checks establish implementation invariants such as evidence binding, state transitions, and deterministic replay. They do not validate astronomical completeness, physical classification, or learned-model performance.

\section{Prospective Evaluation}
A scientifically stronger successor requires an object-grouped, time-forward cohort that records every eligible and rejected object. Before outcomes are inspected, the study must freeze the parent cohort, policy version, external-service snapshots or hashes, image-review labels, review timing, mature outcome definitions, and analysis plan.

For ranking evaluation, SIDEREA should compare the current interpretable ranker, simple supervised baselines, and any shadow representation model using identical time-forward splits. Primary endpoints should be defined at a fixed human-review budget, for example recall among matured target-class objects at top-$K$, together with calibration and route overlap. For the reportability system, gate ablations should be analyzed separately from ranking so that a model cannot receive credit for bypassing safety/evidence checks.

Injection--recovery can characterize sensitivity of the photometric feature path under controlled synthetic perturbations, but it is not a substitute for real-sky completeness. Reviewer reliability should be measured with blinded duplicate review. External-service failure and censoring must remain explicit analysis states rather than being imputed as negative matches.

Until such a cohort exists, learned routes remain shadow-only and the historical campaign remains an operational audit.

\section{Limitations}
The historical record lacks a closed row-level parent cohort, matured labels for all rejected objects, a stable policy over the full campaign, and complete point-in-time external-service snapshots. The current software verification is local and does not constitute live-service qualification. Synthetic transient-search and noise-stress experiments in the repository are useful engineering/research diagnostics but do not establish real-sky completeness or prospective classifier performance. TNS designation is not equivalent to spectroscopic classification. Finally, fixed-radius catalog associations are operational veto/review rules rather than posterior probabilities of physical association.

\section{Conclusion}
SIDEREA is best understood as an evidence and review control plane for optical transient triage. Its central design choice is to prevent a high ranking score from silently becoming a novelty, classification, or reportability claim. The historical campaign provides a bounded operational record, including a large policy-regime discontinuity and two TNS designations, but lacks the denominator required for population-level performance claims. The current implementation adds immutable evidence binding, freshness checks, version-bound review, deterministic finite-budget queues, and shadow-only learned representations. Its source-bound software verification is strong enough to support implementation claims, not astronomical performance claims. A prospective, object-grouped, time-forward cohort is the next required gate for completeness, purity, learned-model comparison, and selection-function inference.

\bibliographystyle{plainnat}
\bibliography{references}

\end{document}
